\section{Modelica Plant Modelling}

\subsection{Plant Components}
PRIMA DI TUTTO IL LAYAOUT

\subsubsection{SMR}

\subsubsection{H2 Plant}

\subsubsection{Methanation Reactor Modeling Approach}

The methanation reactor was modeled in \texttt{Modelica} to simulate the conversion of hydrogen (H\textsubscript{2}) and carbon dioxide (CO\textsubscript{2}) into methane (CH\textsubscript{4}) and water (H\textsubscript{2}O) via the Sabatier reaction:
\begin{equation}
\mathrm{CO_2 + 4H_2 \rightarrow CH_4 + 2H_2O} \qquad \Delta H = -165.1\ \text{kJ/mol}
\end{equation}

To approximate the level of CH\textsubscript{4} produced from a given feed, a dimensionless parameter \texttt{conversionEfficiency} was introduced. This factor (default: 0.78) reflects the maximum chemical efficiency of catalytic methanation as reported in Götz et al.~\cite{Gotz2016}. It represents the fraction of the limiting reactant (either CO\textsubscript{2} or H\textsubscript{2}) that is effectively converted to CH\textsubscript{4}, accounting for reactor inefficiencies and incomplete conversion at realistic temperature and pressure conditions.

The molar output flows are then calculated via stoichiometry:
\begin{align*}
\dot{n}_{\mathrm{CH}_4} &= \dot{n}_{\text{lim}} \cdot \texttt{conversionEfficiency} \\
\dot{n}_{\mathrm{H}_2O} &= 2 \cdot \dot{n}_{\mathrm{CH}_4} \\
\dot{n}_{\mathrm{H}_2} &= \dot{n}_{\mathrm{H}_2,\text{in}} - 4 \cdot \dot{n}_{\mathrm{CH}_4} \\
\dot{n}_{\mathrm{CO}_2} &= \dot{n}_{\mathrm{CO}_2,\text{in}} - \dot{n}_{\mathrm{CH}_4}
\end{align*}

\subsubsection{Turbine}

\subsubsection{Other BOP Components}
Control room, buffers etc etc


\subsection{Demand Definition}
To run the model we decided to compute a priori the daily setpoint for:
\begin{itemize}
  \item Turbine load
  \item HTSE load
  \item Steam tap flow
\end{itemize} 
to do so ..
SPIEGA PROCEDIMENTO DA INPUT FOR MODELICA SECTION DI ENERGY_DEMAND
\subsubsection{Dati che diamo in input}
\subsubsection{Rate limiting}

\subsubsection{Clusters}
spiega che famo i clusters + cita che usiamo dei profili di alcuni gioni particolari per vedere rispospta del sistema